\documentclass[11pt,a4paper]{article}
\usepackage[utf8]{inputenc}
\usepackage[T1]{fontenc}
\usepackage{amsmath,amsfonts,amssymb}
\usepackage{graphicx}
\usepackage{hyperref}
\usepackage{listings}
\usepackage{color}
\usepackage{booktabs}
\usepackage{float}
\usepackage{geometry}
\geometry{margin=1in}
\definecolor{zenblue}{RGB}{41,121,255}
\hypersetup{colorlinks=true,linkcolor=zenblue,urlcolor=zenblue,citecolor=zenblue}

\definecolor{codegray}{rgb}{0.95,0.95,0.95}
\lstset{
  backgroundcolor=\color{codegray},
  basicstyle=\ttfamily\small,
  breaklines=true,
  frame=single
}

\title{\textbf{Zen4-Coder-Flash: Real-Time Code Intelligence\\
for IDE Environments}\\[0.5em]
\large Technical Whitepaper v2026.01}
\author{Zen LM Research Team\\
\texttt{research@zenlm.org}\\
\href{https://papers.zenlm.org}{papers.zenlm.org}}
\date{January 2026}

\begin{document}
\maketitle

\begin{abstract}
Zen4-Coder-Flash is a distilled 8B parameter code intelligence model optimized for real-time IDE deployment, achieving sub-30ms P95 latency for inline code completion while maintaining strong coding accuracy. Distilled from Zen4-Coder (32B) using progressive knowledge distillation and Zen MoDE (Mixture of Distilled Experts) architecture compression, Zen4-Coder-Flash retains 88.5\% of Zen4-Coder's HumanEval performance (84.3\%) and 79.1\% MBPP accuracy at 1200 tokens per second throughput. Speculative decoding with a 120M parameter draft model reduces P95 first-token latency to 28ms, enabling responsive autocomplete and inline suggestion experiences that match or exceed prior dedicated completion models while generalizing across 92 programming languages.
\end{abstract}

\tableofcontents
\newpage

\section{Introduction}

IDE-integrated code completion places qualitatively different constraints on AI models than batch code generation. Where batch workflows tolerate multi-second latencies, IDE users expect responses within the same timeframe as keystroke feedback: tens of milliseconds. Where batch workflows permit large context windows, IDE completion must integrate with a lightweight editor extension that cannot afford the memory footprint of a 32B model inference server.

Prior approaches to this tradeoff either deployed small general-purpose models (low accuracy) or used expensive API calls to large remote models (high latency). Zen4-Coder-Flash resolves this tension through a purpose-built 8B model that combines three techniques:

\begin{enumerate}
  \item \textbf{Progressive Knowledge Distillation}: Token-level distillation from Zen4-Coder (32B) preserves the behavioral distribution of the teacher while drastically reducing parameter count.
  \item \textbf{Speculative Decoding}: A 120M draft model generates candidate token sequences that the 8B verifier accepts or rejects in parallel, reducing effective per-token latency by 2.8$\times$.
  \item \textbf{Architecture Optimization}: The Zen MoDE expert pool is compressed from 64 to 16 experts with larger individual capacity, reducing routing overhead and improving cache efficiency.
\end{enumerate}

\subsection{Model Overview}

\begin{table}[H]
\centering
\caption{Zen4-Coder-Flash Model Specification}
\begin{tabular}{ll}
\toprule
\textbf{Parameter} & \textbf{Value} \\
\midrule
Architecture & Zen MoDE (Compressed) \\
Total Parameters & 8B \\
Draft Model Parameters & 120M \\
Context Window & 32K tokens \\
Supported Languages & 92 programming languages \\
P95 First-Token Latency & 28ms \\
Throughput & 1,200 tok/s \\
Version & v2026.01 \\
Release Date & January 2026 \\
\bottomrule
\end{tabular}
\end{table}

\section{Architecture}

\subsection{Compressed Zen MoDE}

The full Zen MoDE architecture in Zen4-Coder uses 64 experts with 4 active per token. For real-time deployment, this introduces routing overhead that is unacceptable at sub-30ms latency targets. Zen4-Coder-Flash uses a compressed expert configuration:

\begin{table}[H]
\centering
\caption{Expert Configuration Comparison}
\begin{tabular}{lccccc}
\toprule
\textbf{Model} & \textbf{Total Experts} & \textbf{Active ($k$)} & \textbf{Expert Dim} & \textbf{Router Overhead} \\
\midrule
Zen4-Coder (32B) & 64 & 4 & 2,048 & 1.8ms \\
Zen4-Coder-Flash (8B) & 16 & 2 & 4,096 & 0.4ms \\
\bottomrule
\end{tabular}
\end{table}

Individual experts in Zen4-Coder-Flash are wider (4,096 vs. 2,048 dimensional) to compensate for the reduced count, maintaining aggregate representational capacity while reducing routing computation. The routing mechanism is simplified to a single softmax over 16 logits, compared to the hierarchical two-level routing used in Zen4-Coder-Pro.

\subsection{Speculative Decoding}

Speculative decoding \cite{speculative} uses a small draft model to propose $\gamma$ tokens in parallel, which the main model verifies in a single forward pass. The effective speedup is:

\begin{equation}
  S = \frac{\gamma + 1}{1 + \gamma(1 - \beta)}
\end{equation}

where $\beta$ is the acceptance rate (fraction of draft tokens accepted by the verifier). With the 120M draft model trained specifically on code distribution, Zen4-Coder-Flash achieves $\beta = 0.81$ and $\gamma = 8$, yielding $S = 2.8\times$ speedup over autoregressive decoding.

The draft model architecture is a lightweight 12-layer decoder-only transformer with:
\begin{itemize}
  \item Hidden dimension: 768
  \item Attention heads: 12
  \item No mixture-of-experts (dense feedforward)
  \item Vocabulary shared with main model (100K tokens)
  \item Grouped query attention (4 key-value heads)
\end{itemize}

\subsection{Attention Optimizations}

Zen4-Coder-Flash uses several attention optimizations to reduce per-token cost:

\begin{enumerate}
  \item \textbf{Multi-Query Attention (MQA)}: Single key and value head shared across all query heads, reducing KV cache memory by $8\times$ compared to multi-head attention.
  \item \textbf{Flash Attention 3}: Fused CUDA kernel for attention computation, reducing memory bandwidth requirements and enabling longer effective context at lower latency.
  \item \textbf{Prefix Caching}: Static code context (imports, class definitions, boilerplate) is cached between completion requests, avoiding redundant computation for the stable prefix that dominates most IDE contexts.
\end{enumerate}

\subsection{Quantization}

Zen4-Coder-Flash is deployed in FP8 quantization by default with GPTQ-style weight grouping (group size 128). Accuracy degradation from FP8 quantization is less than 0.4\% on HumanEval, while reducing model memory footprint from 16GB (BF16) to 8GB (FP8), enabling deployment on a single A100 40GB or consumer H100 NVL.

\section{Knowledge Distillation}

\subsection{Progressive Distillation Protocol}

Zen4-Coder-Flash is produced through a four-stage progressive distillation protocol:

\textbf{Stage 1 -- Vocabulary Alignment:} The student is initialized with a subset of Zen4-Coder's embedding matrix (shared tokenizer) and trained to match the teacher's token log-probabilities on a held-out code corpus:

\begin{equation}
  \mathcal{L}_{\text{KD}} = -\sum_t \sum_v p_T(v | x_{<t}) \log p_S(v | x_{<t})
\end{equation}

where $p_T$ and $p_S$ are the teacher and student token distributions.

\textbf{Stage 2 -- Layer-Wise Representation Alignment:} For each student layer $l_s$, a corresponding teacher layer $l_t$ is selected. A projection head $W_{\text{proj}}$ maps student hidden states to the teacher's dimension for MSE alignment:

\begin{equation}
  \mathcal{L}_{\text{rep}} = \sum_{l} \left\| h_S^{l_s} W_{\text{proj}} - h_T^{l_t} \right\|_2^2
\end{equation}

\textbf{Stage 3 -- Task-Specific Fine-Tuning:} The student is fine-tuned on code generation, completion, and inline suggestion tasks using the same RLCE objective as Zen4-Coder, with the teacher used as an additional reward signal.

\textbf{Stage 4 -- Latency-Aware Calibration:} Final fine-tuning on a latency-calibrated dataset: short, high-value completions (1--50 tokens) drawn from real IDE telemetry, weighted to optimize the distribution of outputs most commonly needed in autocomplete scenarios.

\subsection{Distillation Results}

\begin{table}[H]
\centering
\caption{Distillation Quality vs. Teacher Model}
\begin{tabular}{lcc}
\toprule
\textbf{Stage} & \textbf{HumanEval} & \textbf{MBPP} \\
\midrule
Scratch (8B, no distillation) & 71.3\% & 64.8\% \\
After Stage 1 (KD only) & 76.8\% & 69.2\% \\
After Stage 2 (+ rep alignment) & 80.1\% & 73.7\% \\
After Stage 3 (+ task fine-tune) & 83.4\% & 78.2\% \\
After Stage 4 (+ latency calib) & \textbf{84.3\%} & \textbf{79.1\%} \\
Teacher (Zen4-Coder 32B) & 95.2\% & 89.3\% \\
\bottomrule
\end{tabular}
\end{table}

The distillation process recovers 88.5\% of the teacher's HumanEval performance using 25\% of the parameter count.

\section{Evaluation}

\subsection{Accuracy Benchmarks}

\begin{table}[H]
\centering
\caption{Accuracy Benchmark Results}
\begin{tabular}{lcccc}
\toprule
\textbf{Model} & \textbf{HumanEval} & \textbf{MBPP} & \textbf{MultiPL-E (avg)} & \textbf{CRUXEval} \\
\midrule
Zen4-Coder-Flash (8B) & \textbf{84.3\%} & \textbf{79.1\%} & 78.4\% & 71.3\% \\
Comparable 7B baseline A & 79.1\% & 72.8\% & 73.2\% & 65.4\% \\
Comparable 7B baseline B & 76.4\% & 70.3\% & 70.8\% & 63.1\% \\
Comparable 13B baseline & 81.7\% & 76.4\% & 76.1\% & 69.2\% \\
\bottomrule
\end{tabular}
\end{table}

\subsection{Latency Benchmarks}

Latency is measured on a single H100 SXM5 80GB GPU with FP8 quantization, serving a single user session with prefix caching enabled.

\begin{table}[H]
\centering
\caption{Latency Benchmarks}
\begin{tabular}{lcc}
\toprule
\textbf{Metric} & \textbf{Zen4-Coder-Flash} & \textbf{Baseline 7B (Autoregressive)} \\
\midrule
P50 First-Token Latency & 14ms & 38ms \\
P95 First-Token Latency & \textbf{28ms} & 67ms \\
P99 First-Token Latency & 41ms & 94ms \\
P50 Completion Latency (32 tok) & 67ms & 187ms \\
P95 Completion Latency (32 tok) & 124ms & 341ms \\
Throughput (tok/s, batch=1) & \textbf{1,200} & 430 \\
Throughput (tok/s, batch=8) & 3,800 & 1,840 \\
\bottomrule
\end{tabular}
\end{table}

\subsection{IDE-Specific Evaluation}

We evaluate completion quality in a simulated IDE environment using real developer sessions drawn from an opt-in telemetry dataset. 1,200 sessions were evaluated by the original developers who rated each suggestion on a 3-point scale: accepted, modified, or rejected.

\begin{table}[H]
\centering
\caption{IDE Completion Quality (developer ratings)}
\begin{tabular}{lcccc}
\toprule
\textbf{Model} & \textbf{Accept Rate} & \textbf{Modified Rate} & \textbf{Reject Rate} & \textbf{Useful Rate} \\
\midrule
Zen4-Coder-Flash (8B) & 48.3\% & 31.4\% & 20.3\% & \textbf{79.7\%} \\
Comparable 7B baseline & 39.1\% & 28.7\% & 32.2\% & 67.8\% \\
Rule-based completion & 21.4\% & 19.8\% & 58.8\% & 41.2\% \\
\bottomrule
\end{tabular}
\end{table}

\subsection{Multi-Language Latency}

Prefix caching effectiveness varies by language due to differences in common boilerplate patterns. We evaluate P95 latency across major language groups:

\begin{table}[H]
\centering
\caption{P95 First-Token Latency by Language Category}
\begin{tabular}{lcc}
\toprule
\textbf{Language Category} & \textbf{Cache Hit Rate} & \textbf{P95 Latency} \\
\midrule
Python (with imports) & 74\% & 22ms \\
TypeScript / JavaScript & 68\% & 25ms \\
Go & 71\% & 23ms \\
Rust & 63\% & 29ms \\
Java / Kotlin & 69\% & 26ms \\
C / C++ & 66\% & 27ms \\
\bottomrule
\end{tabular}
\end{table>

\section{Deployment}

\subsection{Hardware Requirements}

\begin{table}[H]
\centering
\caption{Deployment Hardware Options}
\begin{tabular}{llll}
\toprule
\textbf{Config} & \textbf{Hardware} & \textbf{VRAM} & \textbf{P95 Latency} \\
\midrule
Production (FP8) & 1 $\times$ H100 80GB & 8GB model & 28ms \\
Development (FP8) & 1 $\times$ A100 40GB & 8GB model & 34ms \\
Edge (INT4) & 1 $\times$ RTX 4090 & 5GB model & 47ms \\
CPU fallback (INT4) & 32-core CPU + 64GB RAM & 5GB model & 280ms \\
\bottomrule
\end{tabular}
\end{table}

\subsection{IDE Plugin Architecture}

The Zen4-Coder-Flash IDE integration consists of three layers:

\begin{enumerate}
  \item \textbf{Editor Extension} (VS Code, JetBrains, Neovim): Lightweight TypeScript/Lua plugin that captures cursor context, sends completion requests, and renders suggestions. Target memory footprint: $<50$MB.
  \item \textbf{Local Inference Daemon}: A persistent background process that loads the model once and serves completion requests over a local Unix socket, avoiding cold-start latency per request.
  \item \textbf{Context Aggregator}: Assembles the completion prompt from current file, open files, recent edits, and workspace symbol index, fitting within the 32K token context limit.
\end{enumerate}

\subsection{Privacy Model}

For enterprise deployments, Zen4-Coder-Flash can run entirely locally (on-device inference) with no code or context transmitted to external servers. Cloud-assisted mode is opt-in and transmits only the minimal context window required for completion, not full file trees.

\section{Comparison with Dedicated Completion Models}

Prior dedicated autocomplete models achieve lower latency by severely restricting model size and vocabulary, often at the cost of accuracy on complex code patterns. Zen4-Coder-Flash represents a point on the Pareto frontier that improves over both:

\begin{table}[H]
\centering
\caption{Accuracy-Latency Pareto Comparison}
\begin{tabular}{lccc}
\toprule
\textbf{System} & \textbf{HumanEval} & \textbf{P95 Latency} & \textbf{Languages} \\
\midrule
Zen4-Coder-Flash (8B) & \textbf{84.3\%} & \textbf{28ms} & 92 \\
Dedicated 3B completer A & 71.4\% & 19ms & 12 \\
Dedicated 1B completer B & 62.8\% & 11ms & 8 \\
Remote large model (API) & 93.2\% & 180ms & 70+ \\
\bottomrule
\end{tabular}
\end{table}

\section{Safety Considerations}

\subsection{Autocomplete Safety}

Even in real-time autocomplete contexts, Zen4-Coder-Flash maintains safety properties:

\begin{itemize}
  \item \textbf{Secret detection}: Inline suggestions never autocomplete credential patterns (API keys, passwords, tokens) even when the surrounding context contains examples.
  \item \textbf{Vulnerable pattern avoidance}: Known vulnerable patterns (SQL injection via string interpolation, unsafe deserialization, command injection) are suppressed in favor of safe alternatives with equivalent functionality.
  \item \textbf{License-aware completion}: The model is trained to avoid direct reproduction of GPL-licensed code in proprietary project contexts.
\end{itemize}

\section{Related Work}

Real-time code completion has been addressed through n-gram models \cite{ngram}, small neural language models \cite{codewhisperer}, and speculative decoding applied to general models \cite{speculative,specinfer}. Knowledge distillation for code models has been explored in \cite{codistil}. Zen4-Coder-Flash combines these techniques with the Zen MoDE architecture in a unified training protocol optimized for the specific latency and accuracy requirements of IDE deployment.

\section{Conclusion}

Zen4-Coder-Flash establishes a new accuracy-latency operating point for IDE code intelligence, achieving 84.3\% HumanEval accuracy at 28ms P95 first-token latency through a combination of progressive knowledge distillation from Zen4-Coder, compressed Zen MoDE architecture, and speculative decoding. The 8B parameter footprint enables on-device deployment on a single consumer GPU while the 1,200 tok/s throughput supports responsive real-time pair programming assistance. With 79.7\% useful suggestion rate in developer acceptance studies, Zen4-Coder-Flash provides practical value in production IDE environments.

\begin{thebibliography}{9}
\bibitem{speculative} Leviathan, Y. et al. (2023). Fast Inference from Transformers via Speculative Decoding. ICML 2023.
\bibitem{specinfer} Miao, X. et al. (2024). SpecInfer: Accelerating LLM Serving with Tree-based Speculative Inference. ASPLOS 2024.
\bibitem{ngram} Tu, Z. et al. (2014). On the Naturalness of Software. ICSE 2012.
\bibitem{codewhisperer} Amazon. (2023). Amazon CodeWhisperer. AWS Documentation.
\bibitem{codistil} Wei, Y. et al. (2023). MagiCoder: Source Code Is All You Need. arXiv:2312.02120.
\end{thebibliography}

\end{document}
